% GENERATED FILE. Edit canonical lessons, then run npm run generate:books.
% canonical-source-hash: fnv1a64:535713b3b7cb58ca
% canonical-lessons: LA-C219-fama, LA-C219-nuntius, LA-C219-spectaculum, LA-C219-histrio, LA-C219-fabula

\chapter{Reputation, Messenger, Show, Actor, Story}
\label{ch:la-reputation-messenger-show-actor-story}
\hlchaptermodality{219}

\begin{chapteropening}
\textbf{By the end of this chapter:} \emph{I can say reputation, a messenger, a show, an actor and a story.}

Complete the last lesson of chapter 219: I can say reputation, a messenger, a show, an actor and a story.
\end{chapteropening}

\section[fāma, fāmae]{fāma, fāmae — reputation, rumour}

\label{lesson:LA-C219-fama}



\begin{quote}
Before the new one: say the Latin for a punishment, then the Latin for a charge.
\end{quote}

\subsection*{You'll want to know: fāma, fāmae}
\textbf{fāma, fāmae} — \textquotedblleft{}reputation, rumour\textquotedblright{}.

\begin{cousinweb}
\textbf{fāma} is inherited, usually traced to Proto-Indo-European \emph{*bheh2meh2}, \textquotedblleft{}talk, report\textquotedblright{}, from \emph{*bheh2-}, \textquotedblleft{}to speak\textquotedblright{} (EtymDB 2.1); \emph{fārī}, \textquotedblleft{}to speak\textquotedblright{}, has the same root (Lewis \& Short). Greek \emph{phēmē}, \textquotedblleft{}a saying\textquotedblright{}, is a cousin.
\end{cousinweb}

\begin{grammarlens}[title={What you've built}]
One more word for this chapter.
\end{grammarlens}

\subsection*{Your turn}
\begin{itemize}
\raggedright
  \item \emph{Say it:} \emph{fāma, fāmae}
  \item \emph{Say it:} \emph{fāma, fāmae}, once more
  \item \emph{Say it:} say \emph{fāma, fāmae}
  \item \emph{Recall:} say the Latin for a punishment, then the Latin for a charge, then say \emph{fāma, fāmae} again
\end{itemize}

\begin{tcolorbox}[breakable,colback=teal!4,colframe=teal!35!black,title={Before you move on}]
What does fāma, fāmae mean? (\textquotedblleft{}Reputation, rumour\textquotedblright{}.) Say it once more.
\end{tcolorbox}
\section[nūntius, nūntiī]{nūntius, nūntiī — a messenger; news}

\label{lesson:LA-C219-nuntius}



\begin{quote}
Before the new one: say the Latin for a charge, then the Latin for reputation.
\end{quote}

\subsection*{You'll want to know: nūntius, nūntiī}
\textbf{nūntius, nūntiī} — \textquotedblleft{}a messenger; news\textquotedblright{}.

\begin{cousinweb}
\textbf{nūntius} has no sure origin; it is often explained as \emph{*noventius}, \textquotedblleft{}one who brings something new\textquotedblright{}, from \emph{novus}, \textquotedblleft{}new\textquotedblright{} (Lewis \& Short), but that is uncertain.
\end{cousinweb}

\begin{grammarlens}[title={What you've built}]
One more word for this chapter.
\end{grammarlens}

\subsection*{Your turn}
\begin{itemize}
\raggedright
  \item \emph{Say it:} \emph{nūntius, nūntiī}
  \item \emph{Say it:} \emph{nūntius, nūntiī}, once more
  \item \emph{Say it:} say \emph{nūntius, nūntiī}
  \item \emph{Recall:} say the Latin for a charge, then the Latin for reputation, then say \emph{nūntius, nūntiī} again
\end{itemize}

\begin{tcolorbox}[breakable,colback=teal!4,colframe=teal!35!black,title={Before you move on}]
What does nūntius, nūntiī mean? (\textquotedblleft{}A messenger; news\textquotedblright{}.) Say it once more.
\end{tcolorbox}
\section[spectāculum, spectāculī]{spectāculum, spectāculī — a show}

\label{lesson:LA-C219-spectaculum}



\begin{quote}
Before the new one: say the Latin for reputation, then the Latin for a messenger; news.
\end{quote}

\subsection*{You'll want to know: spectāculum, spectāculī}
\textbf{spectāculum, spectāculī} — \textquotedblleft{}a show\textquotedblright{}.

\begin{cousinweb}
\textbf{spectāculum} is built on \emph{spectō}, \textquotedblleft{}to look at, to watch\textquotedblright{} (Lewis \& Short), with \emph{-culum}, an ending for a means or place: something to watch.
\end{cousinweb}

\begin{grammarlens}[title={What you've built}]
One more word for this chapter.
\end{grammarlens}

\subsection*{Your turn}
\begin{itemize}
\raggedright
  \item \emph{Say it:} \emph{spectāculum, spectāculī}
  \item \emph{Say it:} \emph{spectāculum, spectāculī}, once more
  \item \emph{Say it:} say \emph{spectāculum, spectāculī}
  \item \emph{Recall:} say the Latin for reputation, then the Latin for a messenger; news, then say \emph{spectāculum, spectāculī} again
\end{itemize}

\begin{tcolorbox}[breakable,colback=teal!4,colframe=teal!35!black,title={Before you move on}]
What does spectāculum, spectāculī mean? (\textquotedblleft{}A show\textquotedblright{}.) Say it once more.
\end{tcolorbox}
\section[histriō, histriōnis]{histriō, histriōnis — an actor}

\label{lesson:LA-C219-histrio}



\begin{quote}
Before the new one: say the Latin for a messenger; news, then the Latin for a show.
\end{quote}

\subsection*{You'll want to know: histriō, histriōnis}
\textbf{histriō, histriōnis} — \textquotedblleft{}an actor\textquotedblright{}.

\begin{cousinweb}
\textbf{histriō} came from Etruscan: Livy says the Etruscans called a player \emph{hister}, \textquotedblleft{}actor\textquotedblright{} (Lewis \& Short).
\end{cousinweb}

\begin{grammarlens}[title={What you've built}]
One more word for this chapter.
\end{grammarlens}

\subsection*{Your turn}
\begin{itemize}
\raggedright
  \item \emph{Say it:} \emph{histriō, histriōnis}
  \item \emph{Say it:} \emph{histriō, histriōnis}, once more
  \item \emph{Say it:} say \emph{histriō, histriōnis}
  \item \emph{Recall:} say the Latin for a messenger; news, then the Latin for a show, then say \emph{histriō, histriōnis} again
\end{itemize}

\begin{tcolorbox}[breakable,colback=teal!4,colframe=teal!35!black,title={Before you move on}]
What does histriō, histriōnis mean? (\textquotedblleft{}An actor\textquotedblright{}.) Say it once more.
\end{tcolorbox}
\section[fābula, fābulae]{fābula, fābulae — a story, a play}

\label{lesson:LA-C219-fabula}



\begin{quote}
Before the new one: say the Latin for a show, then the Latin for an actor.
\end{quote}

\subsection*{You'll want to know: fābula, fābulae}
\textbf{fābula, fābulae} — \textquotedblleft{}a story, a play\textquotedblright{}.

\begin{cousinweb}
\textbf{fābula} is built on \emph{fārī}, \textquotedblleft{}to speak\textquotedblright{} (Lewis \& Short), with \emph{-bula}, an ending for a means or result: what is told. English \emph{fable} comes from it.
\end{cousinweb}

\begin{grammarlens}[title={What you've built}]
One more word for this chapter.
\end{grammarlens}

\subsection*{Your turn}
\begin{itemize}
\raggedright
  \item \emph{Say it:} \emph{fābula, fābulae}
  \item \emph{Say it:} \emph{fābula, fābulae}, once more
  \item \emph{Say it:} say \emph{fābula, fābulae}
  \item \emph{Recall:} say the Latin for a show, then the Latin for an actor, then say \emph{fābula, fābulae} again
\end{itemize}

\begin{tcolorbox}[breakable,colback=teal!4,colframe=teal!35!black,title={Before you move on}]
What does fābula, fābulae mean? (\textquotedblleft{}A story, a play\textquotedblright{}.) Say it once more.
\end{tcolorbox}
